\section{Sistemas: tareas, entornos y métodos simples y generales}

La sección anterior explicó por qué el lenguaje es la puerta de entrada; esta sección aborda el núcleo del Agent como sistema. Yao Shunyu (姚顺雨) resume su investigación en dos preguntas: primero, cómo construir tareas y entornos valiosos y más vinculados con el mundo real; segundo, cómo construir métodos simples pero generales. La primera es la línea de las tareas; la segunda, la línea de los métodos. Muchas discusiones solo miran la línea de los métodos y descuidan la de las tareas, pero las capacidades de un Agent suelen estar definidas por el entorno.

\begin{figure}[H]
\centering
\includegraphics[width=0.96\textwidth]{agent-two-core-questions.png}
\caption{Los dos ejes de la investigación sobre Agents: tareas y entornos valiosos + métodos simples y generales. Diagrama conceptual propio, elaborado a partir de la conversación entre 00:17:50 y 00:35:00.}
\end{figure}

\begin{knowledgebox}{Lectura de la figura: tareas y métodos deben evolucionar juntos}
Si la tarea es demasiado simple, el método parece muy potente pero carece de relevancia práctica; si la tarea es lo bastante realista pero el método es demasiado complejo, resulta difícil generalizarlo. La investigación sobre Agents debe diseñar a la vez el entorno, la recompensa, las herramientas y los métodos generales.
\end{knowledgebox}

\subsection{Las tres olas de auge y declive de los Agents: la línea de los métodos y la línea de las tareas}

La sección anterior afirmó que tareas y métodos deben evolucionar juntos; esta sección sitúa ese juicio en la historia de los Agents. Agent es un concepto antiguo: cualquier sistema capaz de tomar decisiones por sí mismo, interactuar con un entorno y optimizar una recompensa puede llamarse Agent. La entrevista nos recuerda que la historia de los Agents no es una línea recta, sino una sucesión de auges y declives: los Agents de la IA clásica y del RL, los entornos de aprendizaje por refuerzo y, después, los LLM Agents. Cada ola no fue solo un cambio de métodos; también dependió de que las tareas y los entornos fueran lo bastante adecuados.

\begin{figure}[H]
\centering
\includegraphics[width=0.96\textwidth]{agent-history-waves.png}
\caption{Las tres olas de auge y declive de los Agents: la línea de los métodos y la de las tareas se complementan; no basta con mirar los algoritmos. Diagrama conceptual propio, elaborado a partir de la conversación entre 00:17:50 y 00:29:00.}
\end{figure}

\begin{knowledgebox}{Términos clave: componentes básicos de un sistema de Agent}
\begin{tabularx}{\textwidth}{p{0.20\textwidth}p{0.32\textwidth}X}
\toprule
Componente & Significado & Por qué importa \\
\midrule
Environment & El mundo que el Agent observa y en el que actúa & Determina si la tarea es realista y si la retroalimentación es válida. \\
Policy & La estrategia que elige acciones según las observaciones & Puede implementarse con un LLM, código, un planificador o un sistema de varios modelos. \\
Reward & La señal que evalúa el resultado de las acciones & Determina qué aprende y qué explora el Agent. \\
Memory & Estado de largo plazo y registro de experiencias & Sostiene el trabajo entre tareas, entre sesiones y la personalización. \\
Tools & Capacidades externas como API, código, navegador y archivos & Permiten que el Agent pase del lenguaje a un mundo ejecutable. \\
\bottomrule
\end{tabularx}
\end{knowledgebox}

\begin{knowledgebox}{Línea de los métodos vs línea de las tareas}
\begin{tabularx}{\textwidth}{p{0.22\textwidth}p{0.34\textwidth}X}
\toprule
Línea & En qué se enfoca & Qué suele pasar por alto \\
\midrule
Línea de los métodos & Algoritmos, modelos, planificación, técnicas de entrenamiento & Si la tarea es lo bastante realista y si el entorno puede dar retroalimentación. \\
Línea de las tareas & benchmark, entornos, herramientas, datos y evaluación & Si el método es simple y general, y si puede transferirse a tareas nuevas. \\
Línea del sistema & Cómo métodos y tareas forman un ciclo cerrado & Costo de cómputo, recuperación ante fallos, memoria de largo plazo y seguridad. \\
\bottomrule
\end{tabularx}
\end{knowledgebox}

Esta tabla también explica por qué muchas demos de Agents parecen impresionantes y, sin embargo, difícilmente se convierten en productos estables. Una demo suele mostrar solo la línea de los métodos: el modelo planifica, invoca herramientas y escribe algunos pasos de acción; un producto real también debe resolver la línea de las tareas: si el entorno es reproducible, si la retroalimentación puede evaluarse de forma inequívoca, si los fallos son recuperables y si el usuario está dispuesto a confiarle sus objetivos reales. Cuando Yao Shunyu insiste en las tareas y los entornos, recuerda precisamente que la investigación sobre Agents no puede limitarse a «aparentar que sabe hacer las cosas».

\begin{importantbox}{Criterios de diseño de un benchmark de Agents}
Un buen benchmark de Agents debe cumplir al menos tres condiciones: tareas lo bastante realistas, retroalimentación lo bastante clara y un entorno lo bastante abierto. Las tareas demasiado simples no permiten medir la generalización; los entornos demasiado cerrados no permiten entrenar acciones reales, y las tareas sin retroalimentación confiable no generan una señal de aprendizaje.
\end{importantbox}

\subsection{Code es la affordance clave}

Hasta aquí se discutieron los componentes de un Agent; esta sección aborda la «interfaz de acción». En la entrevista aparece un juicio muy importante: el code es como la mano humana, la affordance más importante de la IA. Affordance designa las posibilidades de acción que el entorno ofrece a quien actúa. Para la IA, el código, las API, la línea de comandos, el navegador y el sistema de archivos permiten que el modelo no solo hable, sino que modifique el mundo digital.

\begin{figure}[H]
\centering
\includegraphics[width=0.96\textwidth]{code-affordance.png}
\caption{Code es la Affordance clave: el código permite que la IA pase del lenguaje a la acción en el mundo digital. Diagrama conceptual propio, elaborado a partir de la conversación entre 00:29:49 y 00:33:00.}
\end{figure}

\begin{importantbox}{Por qué el code es la mano}
El lenguaje expresa intenciones; el code ejecuta acciones. Un modelo que solo maneja lenguaje es como «una persona que sabe pensar y hablar»; solo el Agent que escribe código, invoca herramientas y modifica el entorno empieza a tener manos en el mundo digital.
\end{importantbox}

\begin{knowledgebox}{Términos clave: Affordance}
Affordance proviene de la psicología ecológica y designa las posibilidades de acción que el entorno ofrece a quien actúa. La manija de una puerta ofrece la affordance de «abrir la puerta»; un botón ofrece la affordance de «hacer clic»; para la IA, el código, las API, la línea de comandos y el navegador ofrecen la affordance de «modificar el mundo digital».
\end{knowledgebox}

\begin{knowledgebox}{Niveles de affordance en el mundo digital}
\begin{tabularx}{\textwidth}{p{0.20\textwidth}p{0.34\textwidth}X}
\toprule
Nivel & Ejemplos & Qué puede hacer el Agent \\
\midrule
Nivel de texto & Documentos, páginas web, correos & Leer, resumir, reescribir, buscar. \\
Nivel de código & Python, SQL, shell, API & Calcular, consultar, invocar servicios, modificar el estado del sistema. \\
Nivel de interfaz & Navegador, IDE, sistema operativo & Ejecutar tareas a través de varias herramientas y manejar entornos no estructurados. \\
Nivel organizacional & Colaboración entre varias personas, permisos, aprobaciones & Completar tareas de alto riesgo junto con los procesos humanos. \\
\bottomrule
\end{tabularx}
\end{knowledgebox}

\subsection{Resumen de la sección}

La clave de un sistema de Agent es el diseño conjunto de tareas, entornos, métodos y affordance. El modelo de lenguaje aporta la capacidad cognitiva; el código y las herramientas, la capacidad de acción; el entorno y la reward, la dirección del aprendizaje.
